\section{The WHOLE original rank-one analytic space is its actual native valued-point space}

\begin{theorem}[The original analytic point map is the actual whole original chart inverse]
\label{thm:fc-native-rank-one-analytic-chart-inverse}
\leanok
\lean{EnvelopeProperty.MultiModel.nativeRankOneOriginalRelativeChart, EnvelopeProperty.MultiModel.nativeRankOnePoint_eq_chart_inverse, EnvelopeProperty.MultiModel.nativeRankOnePoint_mem_relativeLocus, EnvelopeProperty.MultiModel.nativeRankOneRelativePoint, EnvelopeProperty.MultiModel.nativeRankOneRelativePoint_onChart}
\uses{thm:fc-native-rank-one-relative-charts, thm:fc-native-rank-one-analytic-charts}
The original affine analytic chart, followed by the actual native
relative chart inverse, is the true projective-stalk residue point of
the SAME original analytic point. Hence this point lies over the SAME
whole original trivial-base valued point.
\end{theorem}
\begin{proof}
Both whole native points lie in the SAME original coordinate open.
The actual whole absolute chart homeomorphism distinguishes its native
points. Every element of its SAME original chart algebra has equal
value at both points by the two proved chart evaluation formulas.
Thus the native points agree. Original normalization supplies such a
chart for every original analytic point, proving its exact original
base membership without an additional comparison input.
\end{proof}

\begin{theorem}[The whole original rank-one native analytic comparison]
\label{thm:fc-native-rank-one-analytic-comparison}
\leanok
\lean{EnvelopeProperty.MultiModel.nativeRankOneRelativePoint_injective, EnvelopeProperty.MultiModel.nativeRankOneRelativePoint_surjective, EnvelopeProperty.MultiModel.continuous_nativeRankOneRelativePoint, EnvelopeProperty.MultiModel.nativeRankOneAnalyticHomeomorph, EnvelopeProperty.MultiModel.nativeRankOneAnalyticHomeomorph_val, EnvelopeProperty.MultiModel.nativeRankOneAnalyticHomeomorph_support, EnvelopeProperty.MultiModel.nativeRankOneAnalyticHomeomorph_onChart}
\uses{thm:fc-native-rank-one-analytic-chart-inverse, thm:fc-native-rank-one-coordinate-cover}
For every original rank-one model, its WHOLE original analytic space is
homeomorphic to the native valued-point space of its SAME original
projective scheme over its SAME original coefficient field with the
trivial valuation. The homeomorphism retains the true original
projective-stalk residue point, whole original prime kernel, every
original coordinate domain and every whole original affine chart.
\end{theorem}
\begin{proof}
Package the proved original projective residue point with its actual
base membership. If two native images agree, choose an original chart
of the first original point. Equality of the true scheme supports puts
the second point in the SAME original chart. Both actual relative
chart maps and the original affine chart are injective, so the
original points agree. Conversely, the genuine M2-derived finite
native coordinate cover puts any native point in an actual original
coordinate chart. A coordinate nonzero at this point is an original
nonzero coordinate, so its true original chart algebra applies. The
inverse original affine chart reconstructs an original analytic point,
and the proved whole chart-inverse identity recovers the SAME given
native point. On each original open chart the point map is the
composition of the original affine homeomorphism and the actual native
relative chart inverse, hence is continuous; these charts cover.
Finally the original analytic space is compact, and the genuine native
valued-point space is Hausdorff because the actual Proj structure map
and actual affine coefficient map are separated. The continuous whole
bijection is therefore a homeomorphism. All support and chart
identities retain the literal original point construction.
\end{proof}
